\documentclass[../../main.tex]{subfiles}
\graphicspath{{../../image/}{../../figure/}}
\begin{document}

\section{作为压缩半群轨迹的解与周期解的大时间行为}\label{section:ch06-压缩半群与周期解}

对 $t\in[0,\infty)$，考虑映射 $S(t)$，它把 $u_{0}\in L^{\infty}(\mathbb{R}^{m})\cap L^{1}(\mathbb{R}^{m})$ 映为\eqref{align:ch06-6-1-1}、\eqref{align:ch06-6-1-2}的可容许弱解 $u$ 在时刻 $t$ 的限制，即 $S(t)u_{0}(\cdot)=u(\cdot,t)$。由前两节所示可容许弱解的性质，$S(t)$ 把 $L^{\infty}(\mathbb{R}^{m})\cap L^{1}(\mathbb{R}^{m})$ 映入 $L^{\infty}(\mathbb{R}^{m})\cap L^{1}(\mathbb{R}^{m})$，且
\begin{align}
S(0)=I\quad(\text{恒等映射}),
\label{align:ch06-6-4-1}
\end{align}
\begin{align}
S(t+\tau)=S(t)S(\tau),\qquad\text{对任意 }t,\tau\in[0,\infty),
\label{align:ch06-6-4-2}
\end{align}
\begin{align}
S(\cdot)u_{0}\in C^{0}([0,\infty);L^{1}(\mathbb{R}^{m})),
\label{align:ch06-6-4-3}
\end{align}
\begin{align}
\|S(t)u_{0}-S(t)\bar{u}_{0}\|_{L^{1}(\mathbb{R}^{m})}\leqslant\|u_{0}-\bar{u}_{0}\|_{L^{1}(\mathbb{R}^{m})},\qquad\text{对任意 }t\in[0,\infty).
\label{align:ch06-6-4-4}
\end{align}
因而 $S(\cdot)$ 是 $L^{\infty}(\mathbb{R}^{m})\cap L^{1}(\mathbb{R}^{m})$ 上的 L$^{1}$ 压缩半群。

自然产生的问题是：能否通过 Banach 空间中非线性压缩半群的理论从头构造 $S(\cdot)$。这将为\eqref{align:ch06-6-1-1}、\eqref{align:ch06-6-1-2}可容许弱解的存在性提供一个直接、独立的证明，并给出其性质的另一推导。

为构造该半群，须把\eqref{align:ch06-6-1-1}实现为抽象微分方程
\begin{align}
\frac{\mathrm{d}u}{\mathrm{d}t}+A(u)\ni 0,
\label{align:ch06-6-4-5}
\end{align}
其中 $A$ 为适当定义的非线性变换，定义域 $\mathscr{D}(A)$ 与值域 $\mathscr{R}(A)$ 在 $L^{1}(\mathbb{R}^{m})$ 中。一般而言该算子可能是多值的，即对每个 $u\in\mathscr{D}(A)$，$A(u)$ 是 $L^{1}(\mathbb{R}^{m})$ 的非空子集，可含不止一点。

对光滑 $u$，应预期 $A(u)=\operatorname{div}G(u)$。然而把 $\mathscr{D}(A)$ 延拓到非光滑 $u$ 绝非易事，因为构造须以某种方式体现\cref{definition:ch06-6-2-1}中编码的可容许性条件。先作一次初步延拓。为方便，归一化 $G$ 使 $G(0)=0$。

\begin{definition}\label{definition:ch06-6-4-1}
（可能多值的）变换 $\hat{A}$，其定义域 $\mathscr{D}(\hat{A})\subset L^{1}(\mathbb{R}^{m})$，定义为：$u\in\mathscr{D}(\hat{A})$、$w\in\hat{A}(u)$，当且仅当 $u$、$w$、$G(u)$ 均属于 $L^{1}(\mathbb{R}^{m})$，且对任一凸熵函数 $\eta$（$\eta^{\prime}$ 在 $\mathbb{R}$ 上有界，相应熵流 $Q$ 由\eqref{align:ch06-6-2-1}确定）以及 $\mathbb{R}^{m}$ 上所有具紧支集的非负 Lipschitz 连续试验函数 $\psi$，不等式
\begin{align}
\int_{\mathbb{R}^{m}}\Bigl{\{}\sum_{\alpha=1}^{m}\partial_{\alpha}\psi(x)\,Q_{\alpha}(u(x))+\psi(x)\eta^{\prime}(u(x))w(x)\Bigr{\}}\,\mathrm{d}x\geqslant 0
\label{align:ch06-6-4-6}
\end{align}
成立。
\end{definition}

对熵--熵流对 $\pm u$、$\pm G(u)$ 应用\eqref{align:ch06-6-4-6}，可验证对任一 $u\in\mathscr{D}(\hat{A})$，在分布意义下
\begin{align}
\hat{A}(u)=\sum_{\alpha=1}^{m}\partial_{\alpha}G_{\alpha}(u)
\label{align:ch06-6-4-7}
\end{align}
成立。特别地，$\hat{A}$ 是单值的。此外，恒等式
\begin{align}
\int_{\mathbb{R}^{m}}\Bigl{\{}\sum_{\alpha=1}^{m}\partial_{\alpha}\psi\,Q_{\alpha}(u)+\psi\eta^{\prime}(u)\sum_{\alpha=1}^{m}\partial_{\alpha}G_{\alpha}(u)\Bigr{\}}\,\mathrm{d}x=0
\label{align:ch06-6-4-8}
\end{align}
对任一 $u\in C_{0}^{1}(\mathbb{R}^{m})$ 及每一熵--熵流对成立，这蕴含 $C_{0}^{1}(\mathbb{R}^{m})\subset\mathscr{D}(\hat{A})$。特别地，$\mathscr{D}(\hat{A})$ 在 $L^{1}(\mathbb{R}^{m})$ 中稠密。对 $u\in C_{0}^{1}(\mathbb{R}^{m})$，$\hat{A}(u)$ 由\eqref{align:ch06-6-4-7}给出。因而 $\hat{A}$ 确实是\eqref{align:ch06-6-4-7}的延拓。

读者或许已注意到\eqref{align:ch06-6-4-6}与\eqref{align:ch06-6-2-3}的相似之处。与\eqref{align:ch06-6-2-3}类似，为验证\eqref{align:ch06-6-4-6}，只需对熵 $\pm u$ 以及\eqref{align:ch06-6-2-5}或\eqref{align:ch06-6-2-6}的熵--熵流对族检验即可。

\begin{definition}\label{definition:ch06-6-4-2}
（可能多值的）变换 $A$，其定义域 $\mathscr{D}(A)\subset L^{1}(\mathbb{R}^{m})$，是 $\hat{A}$ 的图像闭包，即 $u\in\mathscr{D}(A)$、$w\in A(u)$，当且仅当 $(u,w)$ 是 $L^{1}(\mathbb{R}^{m})\times L^{1}(\mathbb{R}^{m})$ 中某序列 $\{(u_{k},w_{k})\}$ 的极限，其中 $u_{k}\in\mathscr{D}(\hat{A})$、$w_{k}\in\hat{A}(u_{k})$。
\end{definition}

下述命题确立 $A$ 的性质，表明它是 $L^{1}(\mathbb{R}^{m})$ 上一个压缩半群的生成元。

\begin{theorem}\label{theorem:ch06-6-4-3}
变换 $A$ 是增殖的，即若 $u$、$\bar{u}\in\mathscr{D}(A)$，则
\begin{align}
\|(u+\lambda w)-(\bar{u}+\lambda\bar{w})\|_{L^{1}(\mathbb{R}^{m})}\geqslant\|u-\bar{u}\|_{L^{1}(\mathbb{R}^{m})},\qquad \lambda>0,\ w\in A(u),\ \bar{w}\in A(\bar{u}).
\label{align:ch06-6-4-9}
\end{align}
\end{theorem}

\begin{proof}
正是增殖性使 $A$ 生成的半群成为压缩的。因而\cref{theorem:ch06-6-4-3}的证明与\cref{theorem:ch06-6-2-3}中 L$^{1}$ 压缩估计\eqref{align:ch06-6-2-9}的论证极为相似。

由\cref{definition:ch06-6-4-2}，只需证明“较小的”变换 $\hat{A}$ 是增殖的。于是取 $u,\bar{u}\in\mathscr{D}(\hat{A})$，令 $w=\hat{A}(u)$、$\bar{w}=\hat{A}(\bar{u})$。考虑 $\mathbb{R}^{m}\times\mathbb{R}^{m}$ 上任一具紧支集的非负 Lipschitz 连续函数 $\phi$。固定 $\bar{x}\in\mathbb{R}^{m}$，对 Kruzkov 族\eqref{align:ch06-6-2-6}的熵--熵流对 $\eta(u;\bar{u}(\bar{x}))$、$Q(u;\bar{u}(\bar{x}))$ 及试验函数 $\psi(x)=\phi(x,\bar{x})$ 写出\eqref{align:ch06-6-4-6}，得
\begin{align}
\int_{\mathbb{R}^{m}}\operatorname{sgn}[u(x)-\bar{u}(\bar{x})]\Bigl{\{}\sum_{\alpha=1}^{m}\partial_{x_{\alpha}}\phi(x,\bar{x})[G_{\alpha}(u(x))-G_{\alpha}(\bar{u}(\bar{x}))]+\phi(x,\bar{x})w(x)\Bigr{\}}\,\mathrm{d}x\geqslant 0.
\label{align:ch06-6-4-10}
\end{align}
交换 $u$ 与 $\bar{u}$ 的地位，对任一固定 $x\in\mathbb{R}^{m}$ 导出\eqref{align:ch06-6-4-10}的类比式：
\begin{align}
\int_{\mathbb{R}^{m}}\operatorname{sgn}[\bar{u}(\bar{x})-u(x)]\Bigl{\{}\sum_{\alpha=1}^{m}\partial_{\bar{x}_{\alpha}}\phi(x,\bar{x})[G_{\alpha}(\bar{u}(\bar{x}))-G_{\alpha}(u(x))]+\phi(x,\bar{x})\bar{w}(\bar{x})\Bigr{\}}\,\mathrm{d}\bar{x}\geqslant 0.
\label{align:ch06-6-4-11}
\end{align}
将\eqref{align:ch06-6-4-10}关于 $\bar{x}$、\eqref{align:ch06-6-4-11}关于 $x$ 在 $\mathbb{R}^{m}$ 上积分，再相加，得
\begin{align}
&\int_{\mathbb{R}^{m}}\int_{\mathbb{R}^{m}}\operatorname{sgn}[u(x)-\bar{u}(\bar{x})]\Bigl{\{}\sum_{\alpha=1}^{m}(\partial_{x_{\alpha}}+\partial_{\bar{x}_{\alpha}})\phi(x,\bar{x})[G_{\alpha}(u(x))-G_{\alpha}(\bar{u}(\bar{x}))]\notag\\
&\qquad+\phi(x,\bar{x})[w(x)-\bar{w}(\bar{x})]\Bigr{\}}\,\mathrm{d}x\,\mathrm{d}\bar{x}\geqslant 0.
\label{align:ch06-6-4-12}
\end{align}
取 $\mathbb{R}$ 上光滑非负、具紧支集、总质量为一的函数 $\rho$（见\eqref{align:ch06-6-2-15}）。取 $\mathbb{R}^{m}$ 上任一具紧支集的非负 Lipschitz 连续试验函数 $\psi$。对充分小的 $\varepsilon>0$，以
\begin{align}
\phi(x,\bar{x})=\varepsilon^{-m}\psi\Bigl{(}\frac{x+\bar{x}}{2}\Bigr{)}\prod_{\beta=1}^{m}\rho\Bigl{(}\frac{x_{\beta}-\bar{x}_{\beta}}{2\varepsilon}\Bigr{)}
\label{align:ch06-6-4-13}
\end{align}
代入\eqref{align:ch06-6-4-12}，令 $\varepsilon\downarrow 0$。注意到
\begin{align}
(\partial_{x_{\alpha}}+\partial_{\bar{x}_{\alpha}})\phi(x,\bar{x})=\varepsilon^{-m}\partial_{\alpha}\psi\Bigl{(}\frac{x+\bar{x}}{2}\Bigr{)}\prod_{\beta=1}^{m}\rho\Bigl{(}\frac{x_{\beta}-\bar{x}_{\beta}}{2\varepsilon}\Bigr{)},
\label{align:ch06-6-4-14}
\end{align}
得
\begin{align}
\int_{\mathbb{R}^{m}}\sigma(x)\Bigl{\{}\sum_{\alpha=1}^{m}\partial_{\alpha}\psi(x)[G_{\alpha}(u(x))-G_{\alpha}(\bar{u}(x))]+\psi(x)[w(x)-\bar{w}(x)]\Bigr{\}}\,\mathrm{d}x\geqslant 0,
\label{align:ch06-6-4-15}
\end{align}
其中 $\sigma$ 是满足下式的某函数：
\begin{align}
\sigma(x)=
\begin{cases}
1, & u(x)>\bar{u}(x),\\
\in[-1,1], & u(x)=\bar{u}(x),\\
-1, & u(x)<\bar{u}(x).
\end{cases}
\label{align:ch06-6-4-16}
\end{align}
取 $\psi$ 使 $|x|<r$ 时 $\psi(x)=1$，$r\leqslant|x|<r+1$ 时 $\psi(x)=1+r-|x|$，$r+1\leqslant|x|<\infty$ 时 $\psi(x)=0$，令 $r\to\infty$，得
\begin{align}
\int_{\mathbb{R}^{m}}\sigma(x)[w(x)-\bar{w}(x)]\,\mathrm{d}x\geqslant 0,
\label{align:ch06-6-4-17}
\end{align}
其中 $\sigma$ 如\eqref{align:ch06-6-4-16}。

现取任意 $\lambda>0$，利用\eqref{align:ch06-6-4-17}、\eqref{align:ch06-6-4-16}，得
\begin{align}
\|(u+\lambda w)-(\bar{u}+\lambda\bar{w})\|_{L^{1}(\mathbb{R}^{m})}
&\geqslant\int_{\mathbb{R}^{m}}\sigma(x)\{u(x)-\bar{u}(x)+\lambda[w(x)-\bar{w}(x)]\}\,\mathrm{d}x\notag\\
&\geqslant\int_{\mathbb{R}^{m}}\sigma(x)[u(x)-\bar{u}(x)]\,\mathrm{d}x=\|u-\bar{u}\|_{L^{1}(\mathbb{R}^{m})}.
\label{align:ch06-6-4-18}
\end{align}
证明完成。
\end{proof}

\cref{theorem:ch06-6-4-3}结论的一个直接推论（实际上是等价的另一陈述形式）是：

\begin{corollary}\label{corollary:ch06-6-4-4}
对任意 $\lambda>0$，$(I+\lambda A)^{-1}$ 是良定义、单值、L$^{1}$ 压缩的变换，定义在 $I+\lambda A$ 的值域 $\mathscr{R}(I+\lambda A)$ 上。
\end{corollary}

\begin{theorem}\label{theorem:ch06-6-4-5}
变换 $A$ 是极大的，即
\begin{align}
\mathscr{R}(I+\lambda A)=L^{1}(\mathbb{R}^{m}),\qquad\text{对任意 }\lambda>0.
\label{align:ch06-6-4-19}
\end{align}
\end{theorem}

\begin{proof}
由\cref{definition:ch06-6-4-2}与\cref{corollary:ch06-6-4-4}，只需证明 $\mathscr{R}(I+\lambda\hat{A})$ 在 $L^{1}(\mathbb{R}^{m})$ 中稠密；例如只需证明它包含 $L^{1}(\mathbb{R}^{m})\cap L^{\infty}(\mathbb{R}^{m})$。于是固定 $f\in L^{1}(\mathbb{R}^{m})\cap L^{\infty}(\mathbb{R}^{m})$，求方程
\begin{align}
u+\lambda\hat{A}(u)=f
\label{align:ch06-6-4-20}
\end{align}
的解 $u\in\mathscr{D}(\hat{A})$。

回顾 $\hat{A}(u)$ 在分布意义下有形如\eqref{align:ch06-6-4-7}的表示。因而求解\eqref{align:ch06-6-4-20}等价于确定一个一阶拟线性偏微分方程的可容许弱解，即\eqref{align:ch06-6-1-1}的稳态类比。

受\cref{section:ch06-黏性消失法}所讨论的黏性消失法启发，我们把\eqref{align:ch06-6-4-20}的解构造为如下椭圆方程组的解当 $\mu\downarrow 0$ 时的极限：

\begin{align}
u(x)+\lambda\operatorname{div}G(u(x))-\mu\Delta u(x)=f(x),\qquad x\in\mathbb{R}^{m}.
\label{align:ch06-6-4-21}
\end{align}
对任一固定 $\mu>0$，\eqref{align:ch06-6-4-21}在 $H^{2}(\mathbb{R}^{m})$ 中存在解。须证明当 $\mu\downarrow 0$ 时，\eqref{align:ch06-6-4-21}的解族一致有界且几乎处处收敛于某函数 $u$，它即\eqref{align:ch06-6-4-20}的解。证明分以下几步。

\begin{lemma}\label{lemma:ch06-6-4-6}
设 $u_{\mu}$、$\bar{u}_{\mu}$ 是\eqref{align:ch06-6-4-21}以 $f$、$\bar{f}\in L^{1}(\mathbb{R}^{m})$ 为相应右端的解，且取值于紧区间 $[a,b]$。则
\begin{align}
\int_{\mathbb{R}^{m}}[u_{\mu}(x)-\bar{u}_{\mu}(x)]^{+}\,\mathrm{d}x\leqslant\int_{\mathbb{R}^{m}}[f(x)-\bar{f}(x)]^{+}\,\mathrm{d}x,
\label{align:ch06-6-4-22}
\end{align}
\begin{align}
\|u_{\mu}-\bar{u}_{\mu}\|_{L^{1}(\mathbb{R}^{m})}\leqslant\|f-\bar{f}\|_{L^{1}(\mathbb{R}^{m})}.
\label{align:ch06-6-4-23}
\end{align}
此外，若
\begin{align}
f(x)\leqslant\bar{f}(x),\qquad\text{a.e.\ 在 }\mathbb{R}^{m}\text{ 上},
\label{align:ch06-6-4-24}
\end{align}
则
\begin{align}
u_{\mu}(x)\leqslant\bar{u}_{\mu}(x),\qquad\text{在 }\mathbb{R}^{m}\text{ 上}.
\label{align:ch06-6-4-25}
\end{align}
特别地，$u$ 与 $\bar{u}$ 的值域均含于 $[a,b]$。
\end{lemma}

\begin{proof}
与\cref{theorem:ch06-6-3-2}的证明非常相似，留待读者完成。
\end{proof}

\begin{lemma}\label{lemma:ch06-6-4-7}
设 $u_{\mu}$ 为\eqref{align:ch06-6-4-21}以 $f\in L^{\infty}(\mathbb{R}^{m})\cap L^{1}(\mathbb{R}^{m})$ 为右端的解。则当 $\mu\downarrow 0$ 时，$\{u_{\mu}\}$ 一致有界且几乎处处收敛于\eqref{align:ch06-6-4-20}的解 $u$。
\end{lemma}

\begin{proof}
对任意 $y\in\mathbb{R}^{m}$，由 $\bar{u}_{\mu}(x)=u_{\mu}(x+y)$ 定义的函数 $\bar{u}_{\mu}$ 是\eqref{align:ch06-6-4-21}以 $\bar{f}(x)=f(x+y)$ 为右端的解。于是由\eqref{align:ch06-6-4-23}，
\begin{align}
\int_{\mathbb{R}^{m}}|u_{\mu}(x+y)-u_{\mu}(x)|\,\mathrm{d}x\leqslant\int_{\mathbb{R}^{m}}|f(x+y)-f(x)|\,\mathrm{d}x.
\label{align:ch06-6-4-26}
\end{align}
因而族 $\{u_{\mu}\}$ 在 $L^{1}$ 中一致有界且一致等度连续。于是每个序列 $\{\mu_{k}\}$（$\mu_{k}\to 0$，$k\to\infty$）都含子列（仍记为 $\{\mu_{k}\}$），使
\begin{align}
u_{\mu_{k}}\to u,\qquad\text{在 }\mathbb{R}^{m}\text{ 上一致有界且 a.e.},
\label{align:ch06-6-4-27}
\end{align}
其中 $u\in L^{\infty}(\mathbb{R}^{m})\cap L^{1}(\mathbb{R}^{m})$。

现考虑任一光滑凸熵函数 $\eta$，相应熵流 $Q$ 由\eqref{align:ch06-6-2-1}确定。则 $u_{\mu}$ 满足恒等式
\begin{align}
\eta^{\prime}(u_{\mu})u_{\mu}+\lambda\operatorname{div}Q(u_{\mu})-\mu\Delta\eta(u_{\mu})+\mu\eta^{\prime\prime}(u_{\mu})|\nabla u_{\mu}|^{2}=\eta^{\prime}(u_{\mu})f.
\label{align:ch06-6-4-28}
\end{align}
以 $\mathbb{R}^{m}$ 上任一具紧支集的非负光滑试验函数 $\psi$ 乘\eqref{align:ch06-6-4-28}并在 $\mathbb{R}^{m}$ 上积分，得
\begin{align}
\int_{\mathbb{R}^{m}}\Bigl{\{}\lambda\sum_{\alpha=1}^{m}\partial_{\alpha}\psi\,Q_{\alpha}(u_{\mu})+\psi\eta^{\prime}(u_{\mu})(f-u_{\mu})\Bigr{\}}\,\mathrm{d}x\geqslant-\mu\int_{\mathbb{R}^{m}}\Delta\psi\,\eta\,\mathrm{d}x.
\label{align:ch06-6-4-29}
\end{align}
由\eqref{align:ch06-6-4-27}与\eqref{align:ch06-6-4-29}，
\begin{align}
\int_{\mathbb{R}^{m}}\Bigl{\{}\sum_{\alpha=1}^{m}\partial_{\alpha}\psi\,Q_{\alpha}(u)+\psi\eta^{\prime}(u)\lambda^{-1}(f-u)\Bigr{\}}\,\mathrm{d}x\geqslant 0,
\label{align:ch06-6-4-30}
\end{align}
这表明 $u$ 确为\eqref{align:ch06-6-4-20}的解。

由\cref{corollary:ch06-6-4-4}，\eqref{align:ch06-6-4-20}的解唯一，故整个族 $\{u_{\mu}\}$ 当 $\mu\downarrow 0$ 时收敛于 $u$。证明完成。
\end{proof}
\end{proof}

一旦确立增殖性与极大性，非自反 Banach 空间中的 Crandall--Liggett 半群理论保证 $A$ 在 $\overline{\mathscr{D}(A)}=L^{1}(\mathbb{R}^{m})$ 上生成压缩半群 $S(\cdot)$。$S(\cdot)u_{0}$ 可通过隐式差分格式求解微分方程\eqref{align:ch06-6-4-5}来构造：
\begin{align}
\begin{cases}
\dfrac{1}{\varepsilon}[u_{\varepsilon}(t)-u_{\varepsilon}(t-\varepsilon)]+A(u_{\varepsilon}(t))\ni 0, & t>0,\\[6pt]
u_{\varepsilon}(t)=u_{0}, & t<0.
\end{cases}
\label{align:ch06-6-4-31}
\end{align}
对任意 $\varepsilon>0$，由\cref{theorem:ch06-6-4-5}与\cref{corollary:ch06-6-4-4}，\eqref{align:ch06-6-4-31}在 $[0,\infty)$ 上存在唯一解 $u_{\varepsilon}$。进而可证，\cref{corollary:ch06-6-4-4}提供了必要的稳定性，使得当 $\varepsilon\downarrow 0$ 时 $u_{\varepsilon}(\cdot)$ 在 $[0,\infty)$ 的紧子集上一致收敛于某函数，记为 $S(\cdot)u_{0}$。

$S(\cdot)$ 的一般性质由 Crandall--Liggett 理论给出：当 $u_{0}\in\mathscr{D}(A)$ 时，$S(t)u_{0}$ 对一切 $t\in[0,\infty)$ 都留在 $\mathscr{D}(A)$ 中。一般而言，即使 $u_{0}\in\mathscr{D}(A)$，$S(t)u_{0}$ 关于 $t$ 也可能不可微。因而 $S(\cdot)u_{0}$ 应理解为微分方程\eqref{align:ch06-6-4-5}的弱解。

$S(\cdot)$ 的特殊性质是由上述命题（如\cref{lemma:ch06-6-4-6}）所诱导的 $A$ 的特殊性质的推论。下述定理的证明见\cref{section:ch06-注记}所引文献，它总结了 $S(\cdot)$ 的性质，并特别地为\eqref{align:ch06-6-1-1}、\eqref{align:ch06-6-1-2}唯一可容许弱解的存在性（\cref{theorem:ch06-6-2-2}）及其基本性质（\cref{theorem:ch06-6-2-3}与\cref{theorem:ch06-6-2-7}）提供了另一证明。


\begin{theorem}\label{theorem:ch06-6-4-8}
变换 $A$ 在 $L^{1}(\mathbb{R}^{m})$ 上生成压缩半群 $S(\cdot)$，即一族映射 $S(t)\colon L^{1}(\mathbb{R}^{m})\to L^{1}(\mathbb{R}^{m})$（$t\in[0,\infty)$），满足半群性质\eqref{align:ch06-6-4-1}、\eqref{align:ch06-6-4-2}；连续性性质\eqref{align:ch06-6-4-3}（对任意 $u_{0}\in L^{1}(\mathbb{R}^{m})$）；以及压缩性质\eqref{align:ch06-6-4-4}（对任意 $u_{0},\bar{u}_{0}\in L^{1}(\mathbb{R}^{m})$）。若
\begin{align}
u_{0}\leqslant\bar{u}_{0},\qquad\text{a.e.\ 在 }\mathbb{R}^{m}\text{ 上},
\label{align:ch06-6-4-32}
\end{align}
则
\begin{align}
S(t)u_{0}\leqslant S(t)\bar{u}_{0},\qquad\text{a.e.\ 在 }\mathbb{R}^{m}\text{ 上}.
\label{align:ch06-6-4-33}
\end{align}
对 $1\leqslant p\leqslant\infty$，集合 $L^{p}(\mathbb{R}^{m})\cap L^{1}(\mathbb{R}^{m})$ 在 $S(t)$ 下正不变，且对任意 $t\in[0,\infty)$，
\begin{align}
\|S(t)u_{0}\|_{L^{p}(\mathbb{R}^{m})}\leqslant\|u_{0}\|_{L^{p}(\mathbb{R}^{m})},\qquad\text{对一切 }u_{0}\in L^{p}(\mathbb{R}^{m})\cap L^{1}(\mathbb{R}^{m}).
\label{align:ch06-6-4-34}
\end{align}
若 $u_{0}\in L^{\infty}(\mathbb{R}^{m})\cap L^{1}(\mathbb{R}^{m})$，则 $S(\cdot)u_{0}$ 是\cref{definition:ch06-6-2-1}意义下\eqref{align:ch06-6-1-1}、\eqref{align:ch06-6-1-2}的可容许弱解。
\end{theorem}

读者应注意，经由半群的途径为\eqref{align:ch06-6-1-1}、\eqref{align:ch06-6-1-2}提示了一种可容许弱解概念，适用于 $L^{1}(\mathbb{R}^{m})$ 中任意甚至无界的 $u_{0}$。除非流通量 $G$ 在无穷远处呈线性增长，这些解不必是\eqref{align:ch06-6-1-1}的分布解。

Banach 空间中压缩半群的理论为描述当初值为周期时 Cauchy 问题\eqref{align:ch06-6-1-1}、\eqref{align:ch06-6-1-2}解的大时间行为提供了适当框架，例如初值满足
\begin{align}
u_{0}(x+e_{i})=u_{0}(x),\qquad x\in\mathbb{R}^{m},\ i=1,\cdots,m,
\label{align:ch06-6-4-35}
\end{align}
其中 $(e_{1},\cdots,e_{m})$ 是 $\mathbb{R}^{m}$ 的标准基。空间周期性对任一固定 $t\in[0,\infty)$ 传递给解 $u$。因而不妨把 $u_{0}(\cdot)$ 与 $u(\cdot,t)$ 实现为标准环面 $\mathbb{T}^{m}$ 上的函数；此时由\eqref{align:ch06-6-2-9}，把 $u_{0}(\cdot)$ 映为 $u(\cdot,t)$ 的映射族 $S(t)$（$t\geqslant 0$）构成 $L^{\infty}(\mathbb{T}^{m})$ 上的 L$^{1}$ 压缩半群。

L$^{1}$ 压缩性质蕴含：半群从 $u_{0}$ 出发的轨迹 $Su_{0}$ 的轨道 $\gamma(u_{0})=\bigcup_{t\geqslant 0}S(t)u_{0}$ 在 $L^{1}(\mathbb{T}^{m})$ 中相对紧。此外，由半群性质，对一切 $t\in[0,\infty)$ 有 $\gamma(S(t)u_{0})\subset\gamma(u_{0})$。因而 $u_{0}$ 的非空紧 omega 极限集
\begin{align}
\omega(u_{0})=\bigcap_{t\geqslant 0}\overline{\gamma(S(t)u_{0})},
\label{align:ch06-6-4-36}
\end{align}
编码了\eqref{align:ch06-6-1-1}、\eqref{align:ch06-6-1-2}解的大时间行为，因为当 $t\to\infty$ 时，对任意 $1\leqslant p<\infty$，$u(\cdot,t)\to\omega(u_{0})$（在 $L^{p}(\mathbb{T}^{m})$ 中）。

omega 极限集是不变的：$S(t)\omega(u_{0})=\omega(u_{0})$（$0\leqslant t<\infty$）；并且是极小的：对任一 $v_{0}\in\omega(u_{0})$，$\omega(u_{0})=\gamma(v_{0})=\omega(v_{0})$。为明了此点，设 $v_{0}=\lim_{n\to\infty}S(t_{n})u_{0}$、$w_{0}=\lim_{k\to\infty}S(\tau_{k})u_{0}$，其中 $t_{n}\to\infty$、$\tau_{k}\to\infty$。取 $\{\tau_{k}\}$ 的任一子列 $\{\tau_{k_{n}}\}$，使 $s_{n}=\tau_{k_{n}}-t_{n}>n$。由于
\begin{align}
\|S(s_{n})v_{0}-w_{0}\|_{L^{1}}\leqslant\|S(s_{n})v_{0}-S(s_{n})S(t_{n})u_{0}\|_{L^{1}}+\|S(\tau_{k_{n}})u_{0}-w_{0}\|_{L^{1}},
\label{align:ch06-6-4-37}
\end{align}
可知 $S(s_{n})v_{0}\to w_{0}$（$n\to\infty$），即 $w_{0}\in\omega(v_{0})$。

$\omega(u_{0})$ 的极小性，结合\cref{theorem:ch06-6-2-7}与 L$^{p}$ 范数的弱下半连续性，蕴含对任意 $1\leqslant p\leqslant\infty$，$\omega(u_{0})$ 位于 $L^{p}(\mathbb{T}^{m})$ 中以 $0$ 为球心的一个球面上。同样易见（文献见\cref{section:ch06-注记}），限制在 $\omega(u_{0})$ 上的半群 $S$ 成为一族 L$^{1}$ 等距，可延拓为一个 L$^{1}$ 等距群 $\hat{S}$。此外，对任一 $v_{0}\in\omega(u_{0})$，$\hat{S}(t)v_{0}$ 是取值于 $L^{1}(\mathbb{T}^{m})$ 的殆周期函数。


流通量为线性时的\eqref{align:ch06-6-1-1}情形表明，满足上一段全部约束的 omega 极限集仍可能相当大。然而，流通量中的非线性诱导阻尼，可把 omega 极限集收缩为单点，即等于解的守恒平均值的常值函数。

$G$ 在由非零向量 $\xi\in\mathbb{R}^{m}$ 所标志的空间方向上的线性退化编码于集合
\begin{align}
\mathscr{N}_{\xi}=\{u\in\mathbb{R}:\xi\cdot G^{\prime\prime}(u)=0\}.
\label{align:ch06-6-4-38}
\end{align}
下述定理给出周期解渐近衰减到其平均值的最优条件。

\begin{theorem}\label{theorem:ch06-6-4-9}
设 $u$ 为 Cauchy 问题\eqref{align:ch06-6-1-1}、\eqref{align:ch06-6-1-2}的可容许解，初值 $u_{0}\in L^{\infty}(\mathbb{R}^{m})$ 满足周期条件\eqref{align:ch06-6-4-35}。若平均值
\begin{align}
\int_{\mathbb{T}^{m}}u_{0}(x)\,\mathrm{d}x=\bar{u}
\label{align:ch06-6-4-39}
\end{align}
对任一非零 $\xi\in\mathbb{Z}^{m}$ 都不是集合 $\mathscr{N}_{\xi}$ 的内点，则
\begin{align}
u(\cdot,t)\to\bar{u},\qquad\text{当 }t\to\infty,
\label{align:ch06-6-4-40}
\end{align}
收敛在 $L^{p}(\mathbb{T}^{m})$ 中成立，对每个 $1\leqslant p<\infty$。
\end{theorem}

\begin{proof}[证明梗概]
不失一般性，设 $u_{0}\in BV(\mathbb{T}^{m})$，此时 $\omega(u_{0})\subset BV(\mathbb{T}^{m})$。固定任一 $v_{0}\in\omega(u_{0})$。须证明 $v_{0}=\bar{u}$（在 $\mathbb{T}^{m}$ 上 a.e.）。由于 $v_{0}$ 的平均值为 $\bar{u}$，只需证明 $\mathbb{T}^{m}$ 中满足 $v_{0}(x)>\bar{u}$ 的集合 $\mathscr{W}$ 测度为零，等价地，证明 BV 函数
\begin{align}
\phi(x)=\chi_{\mathscr{W}}(x)-\int_{\mathbb{T}^{m}}\chi_{\mathscr{W}}(y)\,\mathrm{d}y,\qquad x\in\mathbb{T}^{m},
\label{align:ch06-6-4-41}
\end{align}
在 $\mathbb{T}^{m}$ 上 a.e. 为零。由基本积分几何，可通过证明 $\phi$ 在每个余维一的测地超平面 $\mathbb{P}_{\xi,\rho}=\{x:\xi\cdot x=\rho\}$（$\xi\in\mathbb{Z}^{m}\setminus\{0\}$，$\rho\in\mathbb{R}$）上的积分为零来导出 $\phi$ 为零。其证明见\cref{section:ch06-注记}所引文献，它基于 $\omega(u_{0})$ 的极小性，特别地蕴含 $v_{0}\in\omega(v_{0})$。

\end{proof}

\end{document}
